\documentclass[11pt]{article}
\usepackage[margin=1in]{geometry}
\usepackage{amsmath,amssymb,amsthm,mathtools,booktabs,array,microtype}
\usepackage[hidelinks]{hyperref}
\usepackage{xcolor}
\newtheorem{theorem}{Theorem}[section]
\newtheorem{proposition}[theorem]{Proposition}
\newtheorem{corollary}[theorem]{Corollary}
\newtheorem{remark}[theorem]{Remark}
\newcommand{\SU}{\mathrm{SU}(2)}
\newcommand{\dd}{\mathrm{d}}
\title{A Coherent-Staple Good-Collar Theorem\\
\large Uniform small-time conditional curvature after the frustrated-boundary obstruction}
\author{}
\date{20 July 2026}
\begin{document}
\maketitle

\begin{abstract}
Worst-case conditioning of a Yang--Mills link on arbitrary surrounding links
does not preserve small plaquette curvature: five identity staples and one
negative-center staple move the conditional weak-coupling minimum away from
the identity.  We prove a boundary-restricted replacement.  If the $m$
staples lie in a common geodesic ball of radius $\eta$ in a compact group with
bi-invariant metric, then the conditional heat-kernel measure has, for every
target plaquette outside curvature radius $\rho$, a small-time rate gap at
least $\rho^2-m\eta^2$.  Thus $\eta<\rho/\sqrt m$ is a uniform sufficient
condition.  The condition is gauge covariant and excludes the center-frustrated
boundary, whose minimum common-ball radius is $\pi/2$ in $\SU$.  The result
supplies the good-collar half of a two-tier activity argument; the activity of
collar failure remains a separate obligation.
\end{abstract}

\tableofcontents

\section{Conditional geometry}

Let $G=\SU$ with bi-invariant geodesic distance $d$ and identity $I$.  Orient
the $m$ plaquettes incident on a selected link so that their conditional
factors are $p_t(US_i)$, where $S_i$ is the corresponding staple.  The
conditional measure is
\begin{equation}
 \gamma_{t,S}(\dd U)=Z_{t,S}^{-1}\prod_{i=1}^{m}p_t(US_i)\,\dd U.
 \label{eq:conditional}
\end{equation}
For $\SU$ with heat generator $\tfrac12\Delta$, the exponential scale is
\begin{equation}
 2t\log p_t(g)\longrightarrow-d(I,g)^2.
 \label{eq:heat-rate}
\end{equation}

\begin{proposition}[Gauge-covariant staple coherence]
Suppose there is $C\in G$ such that
\begin{equation}
 \max_{1\leq i\leq m}d(C,S_i)\leq\eta.
 \label{eq:coherence}
\end{equation}
The least such $\eta$ is invariant under the endpoint gauge action on all
staples.
\end{proposition}
\begin{proof}
The endpoint action is a common left and right multiplication of the staple
family.  A bi-invariant metric preserves all distances under this action and
transports a common center $C$ to the transformed common center.
\end{proof}

Set $V=UC$ and $Q_i=C^{-1}S_i$.  Then $d(I,Q_i)\leq\eta$ and the conditional
rate energy becomes
\begin{equation}
 E_Q(V)=\sum_{i=1}^{m}d(I,VQ_i)^2.
 \label{eq:energy}
\end{equation}

\section{Uniform rate theorem}

For target $j$, write
\[
 B_j(\rho)=\{V:d(I,VQ_j)\geq\rho\}.
\]

\begin{theorem}[Coherent-staple conditional rate]
For every coherent staple family satisfying \eqref{eq:coherence},
\begin{equation}
 \inf_{B_j(\rho)}E_Q-\min_G E_Q\geq \rho^2-m\eta^2.
 \label{eq:gap}
\end{equation}
Consequently, if $\eta<\rho/\sqrt m$, then
\begin{equation}
 \limsup_{t\downarrow0}2t\log\gamma_{t,S}(B_j(\rho))
 \leq-(\rho^2-m\eta^2)<0.
 \label{eq:ldp}
\end{equation}
The same rate upper bound holds for the union of the finitely many target
events.
\end{theorem}
\begin{proof}
At the reference point $V=I$,
\[
 E_Q(I)=\sum_i d(I,Q_i)^2\leq m\eta^2,
\]
so $\min E_Q\leq m\eta^2$.  On $B_j(\rho)$, the $j$th nonnegative summand in
\eqref{eq:energy} is at least $\rho^2$, hence
$\inf_{B_j(\rho)}E_Q\geq\rho^2$.  Subtraction proves \eqref{eq:gap}.
The heat-kernel Laplace principle from \eqref{eq:heat-rate} gives
\eqref{eq:ldp}.  A finite union does not change the exponential rate upper
bound.
\end{proof}

\begin{remark}
The theorem uses a deliberately robust energy comparison.  It does not claim
that $\rho/\sqrt m$ is the optimal coherence threshold.  Its value is that it
is uniform over all noncommuting staple families and needs no coordinate
choice.
\end{remark}

\section{The center-frustrated boundary is excluded}

The geodesic distance between $I$ and $-I$ in $\SU$ is $\pi$.  Any ball
containing both has radius at least $\pi/2$, and the midpoint ball attains this
value.  Thus the Stage 10 family containing both center elements has coherence
radius exactly $\pi/2$.  It cannot enter a collar with ultraviolet
$\eta\ll1$.

For comparison, put all coherent staples in a common one-parameter subgroup
and write their signed angles as $q_i\in[-\eta,\eta]$.  In the convex chart,
\[
 E(\theta)=\sum_i(\theta+q_i)^2,
 \qquad \theta_*=-\frac1m\sum_iq_i.
\]
Enumeration of the $2^m$ vertices of the angle box gives the sharp commuting
bound
\begin{equation}
 \max_i|\theta_*+q_i|\leq\frac{2(m-1)}m\eta.
\end{equation}
For $m=6$ this is $5\eta/3$, smaller than the general metric certificate
$\sqrt6\eta$.

\section{Ultraviolet calibration}

At the frozen Stage 09 parameters and the declared collar radius:
\begin{center}
\begin{tabular}{@{}lr@{}}
\toprule
quantity & value\\
\midrule
$t$ & $3\times10^{-6}$\\
$m$ & $6$\\
$\rho$ & $0.0082$\\
$\eta$ & $0.001$\\
$\rho/\sqrt m$ & $0.00334763598$\\
$\rho^2-m\eta^2$ & $0.00006124$\\
$(\rho^2-m\eta^2)/(2t)$ & $10.2066667$\\
exponential proxy & $3.69233\times10^{-5}$\\
\bottomrule
\end{tabular}
\end{center}
The proxy is a calibration of the asymptotic rate, not a finite-$t$ prefactor
theorem.

\section{Frontier audit}

The result proves that good-collar conditioning is mathematically viable and
that the exact center obstruction is outside its hypothesis.  It does not
estimate how often the collar fails under the interacting link measure.  A
full-measure argument must add a connected activity for incoherent collars,
including center, boundary-holonomy, orbit-atlas, and topological labels.  No
repeated RG, continuum construction, or Hamiltonian mass gap follows from
\eqref{eq:ldp}.  The result is a restricted theorem at claim level C2.

\end{document}

